% Part: history
% Chapter: biographies 
% Section: alan-turing 

\documentclass[../../../include/open-logic-section]{subfiles}

\begin{document}

\olfileid{his}{bio}{tur} 

\olsection{ॲलन ट्यूरिंग}

ॲलन ट्यूरिंग यांचा जन्म २३ जून १९१२ रोजी लंडनमधील मायडा व्हेल येथे झाला. त्यांना सैद्धांतिक संगणकशास्त्राचे जनक मानले जाते. भौतिक विज्ञान आणि गणित यांविषयी त्यांना लहानपणापासूनच रस होता. पण ते शिकत असलेल्या शाळांत साहित्य आणि अभिजात विषयांवर भर असल्याने त्यांच्या आवडीला फार वाव मिळाला नाही. परिणामी, शाळेत त्यांची कामगिरी कमकुवत राहिली आणि अनेक शिक्षकांनी त्यांना तंबी दिली.

\olphoto{turing-alan}{ॲलन ट्यूरिंग}

ट्यूरिंग यांनी केंब्रिजच्या किंग्ज कॉलेजमध्ये गणिताचे पदवीशिक्षण घेतले. गणनीय फलन ही संकल्पना नेमकी परिभाषित करण्यासाठी आणि निर्णय समस्या अनिर्णेय आहे हे सिद्ध करण्यासाठी त्यांनी १९३६ मध्ये पुढे ट्यूरिंग यंत्र म्हणून ओळखली गेलेली संकल्पना विकसित केली. ॲलोन्झो चर्च यांनी स्वतःच्या लॅम्डा कलनाच्या साहाय्याने हा निष्कर्ष आधीच सिद्ध केला होता. तरीही चर्च यांच्या निष्कर्षाचा उल्लेख करून ट्यूरिंग यांचा लेख प्रकाशित झाला. चर्च यांच्या निमंत्रणावरून ट्यूरिंग प्रिन्स्टनला गेले. तेथे ते १९३६ ते १९३८ या काळात राहिले आणि चर्च यांच्या मार्गदर्शनाखाली डॉक्टरेट मिळवली.

तर्कशास्त्रातील रस असूनही भौतिक विज्ञानातील ट्यूरिंग यांची जुनी आवड कायम होती. दुसऱ्या महायुद्धात ब्रिटनच्या ब्लेचली पार्क येथील संकेतभेदन विभागात काम करताना त्यांच्या व्यावहारिक कौशल्याचा उपयोग झाला. जर्मन नौदलाच्या संदेशांमध्ये वापरला जाणारा एनिग्मा संकेत भेदण्यात त्यांची मध्यवर्ती भूमिका होती. संख्याशास्त्र आणि संकेतलेखनातील त्यांचे कौशल्य, तसेच इलेक्ट्रॉनिक यंत्रांचा वापर, यांमुळे चमूला “बॉम्ब” (bombe) नावाचे संकेतभेदक यंत्र तयार करून हा संकेत भेदता आला. त्यांच्या कल्पनांनी जगातील पहिला प्रोग्राम करता येणारा इलेक्ट्रॉनिक संगणक, कोलॉसस, तयार करण्यासही मदत केली; ब्लेचली पार्क येथेच तो जर्मन लोरेन्झ संकेत भेदण्यासाठी वापरला गेला.

ट्यूरिंग समलैंगिक होते. तरीही १९४२ मध्ये त्यांनी ब्लेचली पार्कमधील सहकारी जोन क्लार्क यांना लग्नाची मागणी घातली. नंतर त्यांनी साखरपुडा मोडला आणि आपण समलैंगिक आहोत असे त्यांना सांगितले. ट्यूरिंग यांचे आयुष्यात अनेक पुरुष जोडीदार होते; त्या काळी ब्रिटनमध्ये पुरुषांमधील समलैंगिक संबंध कायद्याने गुन्हा होते. १९५२ मध्ये त्यावेळच्या जोडीदाराच्या एका मित्राने ट्यूरिंग यांच्या घरात चोरी केली. पोलिसांत तक्रार करताना ट्यूरिंग यांनी आपले समलैंगिक संबंध असल्याचे सांगितले; अशा संबंधांना लवकरच कायदेशीर मान्यता मिळेल, असा त्यांचा समज होता. मात्र तसे झाले नाही आणि त्यांच्यावर “गंभीर अश्लीलता” (gross indecency) या तत्कालीन गुन्ह्याचा आरोप ठेवण्यात आला. तुरुंगवासाऐवजी ट्यूरिंग यांनी कामेच्छा कमी करणारे संप्रेरक उपचार निवडले. ८ जून १९५४ रोजी सायनाइडचे अतिप्रमाण सेवन झाल्याने ते मृतावस्थेत आढळले; स्रोताच्या मते ती बहुधा आत्महत्या असावी. २०१३ मध्ये राणी एलिझाबेथ दुसरी यांनी त्यांना राजमाफी दिली.

\begin{reading}
ॲलन ट्यूरिंग यांच्या सविस्तर चरित्रासाठी \citet{Hodges2014} पाहा. त्यांच्या जीवनावर आणि कार्यावर आधारित \emph{Breaking the Code} हे नाटक १९९६ मध्ये डेरेक जेकोबी यांनी ट्यूरिंग यांची भूमिका केलेल्या दूरचित्रवाणी कार्यक्रमाच्या रूपात सादर झाले. अकादमी पुरस्कारासाठी नामांकन मिळालेला \emph{The Imitation Game} हा बेनेडिक्ट कंबरबॅच आणि कीरा नाइटली यांच्या भूमिका असलेला चित्रपटही ट्यूरिंग यांचे जीवन आणि ब्लेचली पार्कमधील काळ यांवर सैलपणे आधारलेला आहे \citep{Imitation2014}.

\citet{Radiolab2012} येथे ट्यूरिंग यांच्या जीवन आणि कार्यावरील अनेक पॉडकास्ट आहेत. बीबीसी होरायझनचा \emph{The Strange Life and Death of
  Dr.~Turing} हा माहितीपट ऑनलाइन पाहता येतो \citep{Sykes1992}. ट्यूरिंग यांच्या जन्मशताब्दीनिमित्त २०१२ मध्ये तयार केलेले कार्यरत लेगो ट्यूरिंग यंत्र दाखवणारा छोटा व्हिडिओ \citep{Theelen2012} येथे आहे.

ट्यूरिंग यंत्र आणि निर्णय समस्येवरील ट्यूरिंग यांचा मूळ लेख \citet{Turing1937} हा आहे.
\end{reading}

\end{document}
